\section{Conclusion}
\label{sec:conclusion}

%% [MOVE 1 — CLOSE THE LOOP] The introduction asked how much of the per-stage
%% decision actually needs learning. Answer it directly.
This paper asked how much of the DWTA per-stage decision a learned policy
actually needs to make. The answer is one bit per weapon per stage. Given the
set of weapons that fire, assigning targets among them is monotone submodular
maximization under a partition matroid, which a many-to-one greedy rule on
exact marginal values solves to within half the stage optimum without
training (Proposition~\ref{prop:auction-half}); and because that guarantee
holds for any firing set, it holds for one produced by all weapons deciding
at once, which is why our simultaneous decoder needs neither a conflict
handler nor a communication layer to match what sequential decoding would
have guaranteed (Remark~\ref{rem:auction-kills-pileon}). What cannot be
delegated is the decision to hold: its value is zero within the stage and
realized only later, and any rule that optimizes the current stage destroys
an arbitrarily small fraction of the optimum on instances where a window
closes (Proposition~\ref{prop:myopic-gap}). Learning that decision alone, and
training it against the objective a fixed local search produces rather than
its own raw output, reduced the mean objective by 14.9\% relative to the
strongest alternative we trained, across twelve held-out configurations
spanning an order of magnitude in scale.

%% [MOVE 2 — SO WHAT] Transferable lessons, phrased as things another
%% practitioner can act on.
Two results transfer past DWTA. The first is a criterion for where to cut a
learned pipeline: give the classical subroutine the part of the decision that
admits a per-stage approximation guarantee and keep for the learner the part
for which no per-stage guarantee exists even in principle. That criterion is
checkable before any training run, and in our case it identified a decision
boundary that a natural end-to-end design had obscured. The second is a
control that improvement-learning pipelines should report and, in our reading
of the literature, generally do not: because a keep-best or
accept-if-improving rule makes any proposal distribution monotonically
non-degrading, the comparison that establishes a learned proposer's value is
against a uniformly random proposer at matched re-simulation budget, not
against the unimproved base. Running that control on our own learned editor
was what led us to delete it; a fixed hill-climbing search proved cheaper
and better (Section~\ref{sec:local-search}).

%% [MOVE 3 — LIMITS] Honest boundaries.
Several limits bound these claims. Results are reported over three seeds, and
while the across-seed spread is small (standard deviation between 0.000 and
0.002, below the gap to every competing method at every configuration), three
seeds is a thin basis for the closest per-configuration calls. A sequential
decoder still edges ahead at two configurations, by 0.005 and 0.004, and POMO
at one, by 0.007. Training exhibits
the early-peak-then-oscillate behaviour common to every method we trained on
this problem, so the reported checkpoint is selected by best held-out
performance during training rather than final step; we state this openly and
report the full evaluation trace rather than the peak alone, but it remains a
selection procedure that a single-run reader cannot verify. The fire
probability is computed from per-target scores as well as a hold score
(Eq.~\eqref{eq:hold-prob}); a head with a single scalar fire-logit per weapon
would be simpler still, and we have not tested it. Finally, the local search's operator is a single participation flip.
Richer operators we tested earlier, including one that reschedules a shot
between stages, did not outperform it, but the space is not exhausted.

%% [MOVE 4 — OPEN QUESTIONS] Questions, not a task list.
The open questions we consider most worth answering follow from those limits.
Does the decomposition criterion above identify useful cut points in other
problems that mix a guaranteed within-epoch allocation with an unguaranteed
cross-epoch commitment, such as inventory with perishable stock or maintenance scheduling under resource windows, or is the clean separability of DWTA
unusual? Under what conditions does a learned proposer beat a random one in
an improvement loop at matched budget; our edit space is one bit wide, and we
would expect the learned proposer's advantage to grow with the operator's
arity, but we have not mapped that curve. And since the assignment rule
supplies a per-stage guarantee while the learned policy supplies none, is there a
trajectory-level guarantee for the composition, or a characterization of the
instances on which the participation decision is provably easy?
